% \section{Overview}
\begin{figure}
\begin{center}
    \begin{align*}
       \expr & \to f\usp\vec{\core} \testEq \core & \text{test}\\
            &\OR \compileDSL{\surface} & \text{DSL compilation} \\
            &\OR\expr \ccons \expr & \text{list cons} \\
            &\OR \nil & \text{empty list} \\
            &\OR \ldots & \\
    \end{align*}
\end{center}
    \caption{Host Language Grammar}
    \label{fig:the-host-language}
\end{figure}

Our approach to property-based testing
uses an embedded domain specific language (eDSL), \ourDSL,
into a \textit{host language}-- 
a traditional programming language that a user wishes to test code in
(our practical implementation uses Haskell as its host language, for example.)
The users write eDSL code in what we refer to as the \textit{surface language}--
this language is designed to be expressive and ergonomic to program in.
The surface language is then compiled to a smaller \textit{core language},
which is then used to actually generate tests.
We use \expr to denote generic host language expressions,
\surface to denote generic surface language expressions,
and \core to denote generic core language expressions.


Figure~\ref{fig:the-host-language}
shows a small snippet of the essential elements in the host language.
\bill{Do we want to show data constructors here?}
The most straightforward of these is
standard \textit{cons-lists},
where we write \ccons for cons and \nil for the empty list.
DSL compilation, $\compileDSL{\surface}$,
and tests, $f\usp\vec{\core} \testEq \core$,
really represent functions we implement in the host language to interface with the eDSL.
DSL compilation, $\compileDSL{\surface}$ translates the surface language expression
$\surface$ into a corresponding core language expression $\core$.
A test $f\usp\vec{\core_1} \testEq \core_2$,
takes a host language function $f$,
a list of (core) symbolic expressions $\vec{\core_1}$ corresponding to
the function's arguments,
and a core symbolic expression $\core_2$ corresponding to the function's output.
The test indicates that $f$ should be run on random host language inputs,
\textit{generated} according to $\vec{\core_1}$,
and then the returned value should be \textit{validated}, or checked for correctness,
according to the output $\core_2$.



